\Subhead{2.1.2}{Плоские акустические волны в вязкой теплопроводящей среде}

Для малых возмущений жидкости около покоя уравнения непрерывности и движения имеют вид
\begin{align}
 &\partial_t\rho'+\rho_0\nabla\cdot\bm v=0,\\
 &\rho_0\partial_t\bm v=-\nabla p'+\eta\Delta\bm v+
 (\zeta+\eta/3)\nabla\nabla\cdot\bm v.
\end{align}
Здесь $\eta$ и $\zeta$ — сдвиговая и объёмная вязкости. Уравнение теплопроводности и уравнение состояния связывают $p',\rho',T'$. Без потерь процесс адиабатический: $p'=c_s^2\rho'$, $c_s^2=(\partial p/\partial\rho)_s$. Исключение скорости даёт $p'_{tt}-c_s^2\Delta p'=0$. В плоской бегущей волне $v_\parallel=p'/(\rho_0c_s)$ и
\begin{equation}
 \overline I=\frac{|p_0|^2}{2\rho_0c_s}=\frac{\rho_0c_s|v_0|^2}{2}.
\end{equation}
Вязкость и теплообмен превращают часть упорядоченной энергии в тепло. В слабопоглощающем длинноволновом приближении $k\simeq\omega/c_s+\ii\alpha$, где
\begin{equation}
 \alpha=\frac{\omega^2}{2\rho_0c_s^3}
 \left[\frac43\eta+\zeta+\kappa_T\left(\frac1{c_v}-\frac1{c_p}\right)\right].
\end{equation}
$\kappa_T$ — теплопроводность, $c_p,c_v$ — теплоёмкости единицы массы. Это коэффициент затухания \emph{амплитуды}: интенсивность убывает как $\e^{-2\alpha x}$. Условие $\alpha\ll\omega/c_s$ позволяет вычислять потери по невозмущённой звуковой волне; близко к релаксационным частотам нужна дисперсионная модель среды. Поперечное движение ньютоновской жидкости подчиняется диффузионному уравнению, а не обычной упругой волне.

% Самостоятельная проверка теплопроводного вклада: rho T ds/dt=kappa_T Delta T,
% ds=(c_v/T)dT-(p_T/rho^2)drho, p'=c_T^2 rho'+p_T T'.
% Исключение T' в первом порядке по потерям и тождество
% T p_T^2/(rho^2 c_v^2 c_s^2)=1/c_v-1/c_p дают тепловой член в alpha.
